% =====================================================================
%  1bat-ud3-15.tex  —  HORA 15: Taller de preparació de la prova
%  (sense preàmbul: s'inclou des de main.tex)
% =====================================================================
\horatitol{Sessió 15 — Taller de preparació de la prova}%
{Conèixer l'estructura de la prova i la rúbrica, assajar un simulacre per blocs i fer un pla d'estudi personal.}

\subsection{Idea clau}

Demà tindrem el repàs ben fresc; avui ens hi preparem \textbf{sense sorpreses}. La
prova de la sessió següent té \textbf{quatre blocs}, un per cada criteri que hem
treballat. Saber \emph{exactament} què es demana i com es corregeix fa que l'examen
deixi de ser imprevisible. Després assajarem amb un \textbf{simulacre}.

\begin{idea}
\textbf{Estructura de la prova} (un bloc per criteri):
\begin{enumerate}[leftmargin=1.6em, itemsep=2pt]
  \item \textbf{Bloc 1 (C1) — Domini i punts de tall:} determinar el domini d'una
        funció i els seus talls amb els eixos.
  \item \textbf{Bloc 2 (C2) — Lineals i quadràtiques:} representar i interpretar una
        recta o una paràbola, amb vèrtex si escau.
  \item \textbf{Bloc 3 (C3) — Funcions a trossos:} representar-ne una marcant punts
        oberts i tancats, i comentar-ne la continuïtat.
  \item \textbf{Bloc 4 (C4) — Modelització i interpretació:} llegir una situació real
        (repartiment, tarifa, benefici\dots) i interpretar-la.
\end{enumerate}
\textbf{Es valora}, a més del resultat: la \textbf{representació gràfica acurada},
els \textbf{punts oberts i tancats} ben marcats i la \textbf{interpretació dins del
context}.
\end{idea}

\subsection{Rúbrica simplificada}

\noindent Així es corregeix cada bloc (nivells d'assoliment):
\begin{center}
\renewcommand{\arraystretch}{1.3}
\begin{tabular}{p{2.4cm} p{10.2cm}}
\toprule
\textbf{Nivell} & \textbf{Què demostra} \\
\midrule
\textbf{AE} (excel·lent) & Resol sense errors, representa amb cura i interpreta el
resultat dins del context. \\
\textbf{AN} (notable) & Resol correctament amb algun error menor; la representació i
la interpretació són essencialment bones. \\
\textbf{AS} (suficient) & Resol el procediment bàsic amb alguns errors; representació
o interpretació incompletes. \\
\textbf{NA} (no assolit) & No identifica el procediment o comet errors que
n'invaliden el resultat. \\
\bottomrule
\end{tabular}
\end{center}

\subsection{Simulacre — un ítem de cada bloc}

\begin{exemple}[com es resol un ítem ben fet, de cap a cap]
\textbf{Ítem tipus (Bloc 2):} representa i interpreta $B(x)=-x^2+6x-5$, el benefici
d'una venda solidària segons el preu $x$. Un \textbf{ítem complet} respon: talls,
vèrtex \emph{i} interpretació.
\begin{itemize}[leftmargin=1.4em, itemsep=1pt]
  \item Talls amb $X$: $-x^2+6x-5=0 \Rightarrow x^2-6x+5=0 \Rightarrow (x-1)(x-5)=0
        \Rightarrow x=1,\ x=5$ (preus d'equilibri).
  \item Vèrtex: $x_V=\dfrac{-6}{2\cdot(-1)}=3$, $B(3)=-9+18-5=4$: preu òptim $3\eur$,
        benefici màxim $4\eur$.
  \item \textbf{Interpretació:} entre $1\eur$ i $5\eur$ hi ha benefici; el millor preu
        és $3\eur$. \emph{Aquesta frase final és la que distingeix un AE d'un AS.}
\end{itemize}
\end{exemple}

\begin{exercicis}
\textbf{Resol el simulacre amb temps controlat. Un ítem per bloc:}
\begin{exlist}
  \item \textbf{(Bloc 1 — C1)} Troba el domini i els punts de tall amb els dos eixos
        de $f(x)=\dfrac{x-4}{x+2}$.
  \item \textbf{(Bloc 2 — C2)} Per a la paràbola $f(x)=x^2-2x-8$: talls amb els eixos,
        vèrtex (màxim o mínim?) i un esbós de la gràfica.
  \item \textbf{(Bloc 3 — C3)} Representa la funció a trossos següent i digues si és
        contínua, marcant els punts oberts i tancats:
        \[
        f(x)=
        \begin{cases}
        x+1 & \text{si } x\le 1,\\[2pt]
        4 & \text{si } x>1.
        \end{cases}
        \]
  \item \textbf{(Bloc 4 — C4)} Un fons de $\num{48000}\eur$ es reparteix entre $x$
        famílies: $f(x)=\dfrac{\num{48000}}{x}$. Calcula què rep cada família si n'hi
        ha $8$ i si n'hi ha $40$, i explica què passa amb l'ajut quan creix el nombre
        de famílies.
\end{exlist}
\end{exercicis}

\subsection{Formulari-resum (suport)}

\begin{idea}
\textbf{Tingues-ho a mà mentre prepares la prova:}
\begin{itemize}[leftmargin=1.4em, itemsep=2pt]
  \item \textbf{Domini:} polinomi $\rightarrow$ tot $\mathbb{R}$; fracció
        $\rightarrow$ denominador $\neq 0$; arrel $\rightarrow$ radicand $\ge 0$.
  \item \textbf{Talls:} amb $Y$, fes $x=0$; amb $X$, fes $f(x)=0$.
  \item \textbf{Recta:} $y=mx+n$ ($m$ pendent, $n$ ordenada a l'origen).
  \item \textbf{Vèrtex:} $x_V=\dfrac{-b}{2a}$, després $y_V=f(x_V)$. $a>0$ mínim;
        $a<0$ màxim.
  \item \textbf{A trossos:} cada tram amb el seu interval; punt tancat $\bullet$
        (inclòs), obert $\circ$ (exclòs); contínua si els trossos s'enganxen.
\end{itemize}
\end{idea}

\noindent\textbf{Pla d'estudi personal.} A partir de la teva graella
d'autoavaluació (sessió 9) i de com t'ha anat el simulacre, anota els \textbf{dos o
tres tipus d'exercici} que has de repassar amb més atenció abans de la prova:
\begin{center}
\renewcommand{\arraystretch}{1.5}
\begin{tabular}{c p{10.5cm}}
\toprule
\textbf{\#} & \textbf{Què he de repassar (i com)} \\
\midrule
1 & \rule{9.8cm}{0.4pt} \\
2 & \rule{9.8cm}{0.4pt} \\
3 & \rule{9.8cm}{0.4pt} \\
\bottomrule
\end{tabular}
\end{center}

% ---------------------------------------------------------------------
\fullsolucions{Sessió 15 — simulacre}
\begin{solucions}
\begin{sollist}
  \item \textbf{(Bloc 1)} Domini: denominador zero a $x+2=0 \Rightarrow x=-2$, per
        tant $\mathrm{Dom}=\mathbb{R}\setminus\{-2\}$. Tall amb $Y$ ($x=0$):
        $f(0)=\dfrac{-4}{2}=-2$, punt $(0,-2)$. Tall amb $X$ (numerador zero):
        $x-4=0 \Rightarrow x=4$, punt $(4,0)$.
  \item \textbf{(Bloc 2)} Tall $Y$: $(0,-8)$. Talls $X$: $x^2-2x-8=0
        \Rightarrow (x-4)(x+2)=0 \Rightarrow x=4,\ x=-2$, punts $(4,0)$ i $(-2,0)$.
        Vèrtex: $x_V=\dfrac{2}{2}=1$, $f(1)=1-2-8=-9$, vèrtex $(1,-9)$, \textbf{mínim}
        (s'obre cap amunt).
  \item \textbf{(Bloc 3)} A $x=1$: el tros esquerre dóna $f(1)=2$ (punt tancat
        $(1,2)$); el tros dret val $4$ (punt obert $(1,4)$). Com que $2\neq 4$, la
        funció \textbf{no és contínua}: hi ha un salt a $x=1$.
  \item \textbf{(Bloc 4)} $f(8)=\dfrac{\num{48000}}{8}=\num{6000}\eur$;
        $f(40)=\dfrac{\num{48000}}{40}=\num{1200}\eur$. Quan creix el nombre de
        famílies, l'ajut per família \textbf{decreix} i s'acosta a $0$ (asímptota
        $y=0$): el fons es diluteix com més gent l'ha de compartir.
\end{sollist}
\end{solucions}
